\subsection{Exact finite dyadic Selmer loci and a cyclic-eight detector}

\begin{statement}{A.193}{Proved}
{A literal local basis and finite separation depths}
Retain \(C:y^2=x^5-54x\), \(J,D,T,j_C\) from A.192,
and normalize \(v_2(2)=1\).  Let \(s\in1+4\Z_2\) be the
unique square root of \(-1176066383\), and put
\[
 R_* =(-86,2s),\qquad E_*=[R_*-\infty],\qquad
 \mathcal M=\varprojlim_m J(\Q_2)/2^mJ(\Q_2).
\tag{A.193.1}
\]
The actual local completion has the direct decomposition
\[
 \mathcal M=\Z_2D\oplus\Z_2E_*\oplus(\Z/2)T.
\tag{A.193.2}
\]
For \(n\geq1\), the finite Kummer class of \(j_C(R)\),
\(R\in C(\Q_2)\), belongs to
\(\operatorname{loc}_2\operatorname{Sel}^{2^n}(J/\Q)\)
exactly in the following cases:
\[
\begin{array}{c|c}
 R&\text{necessary and sufficient condition}\\ \hline
 \infty,(0,0),(3,9),(3,-9)&\text{every }n\\
 v_2(x)=2k-1>0&n\leq k-1\\
 v_2(x)=-2k<0&n\leq k\\
 x\in\Z_2^\times,\ x\ne3&n\leq v_2(x-3).
\end{array}
\tag{A.193.3}
\]
In the last row \(v_2(x-3)\geq3\) automatically.

For the actual cyclic-eight cover \(W\to U\) of
(A.192.2), there is an arithmetic quotient with geometric group
\[
 Q_n\simeq\bigl(\Z_2[\zeta_8]/2^n\bigr)\rtimes C_8,
 \qquad |Q_n|=2^{4n+3},\qquad \operatorname{cl}(Q_n)=4n.
\tag{A.193.4}
\]
It is a quotient of the full geometric metabelian group.
Every locally point-or-cusp-matching section of this arithmetic
quotient in the combined class \([-24]\) from (A.192.3)
obeys (A.193.3) on every selected dyadic cover sheet.
In particular the quotient with geometric group of order
\(2^{19}\) and class sixteen excludes the ordinary dyadic centre
\[
                         t_* =\frac{14641}{14587}
\tag{A.193.5}
\]
in that combined class.  The equivalence in (A.193.3) is the
finite Jacobian Selmer test; the quotient-section assertion is
its necessary consequence with all local curve images retained.
\end{statement}

\begin{proof}
\emph{An elementary certificate for the local basis.}
Set \(L=\Q_2(\theta)\), \(\theta^4=54\).
The defining polynomial is Eisenstein at two, so
\(L\) has residue field \(\F_2\), uniformizer \(\theta\),
and \(v_\theta(2)=4\).  The actual odd-degree Kummer map
proved in A.186 is injective on \(J(\Q_2)/2J(\Q_2)\)
and has values in \(L^\times/L^{\times2}\).
The values on \(D,T,E_*\) are respectively
\[
 d=3-\theta,\qquad t_0=-\theta,\qquad
 r=-86-\theta=t_0\alpha,
 \qquad\alpha=1+\frac{43}{27}\theta^3.
\tag{A.193.6}
\]
The point in (A.193.1) exists because
\((-86)^5-54(-86)=4(-1176066383)\) and
\(-1176066383\equiv1\pmod8\).
The square root with residue one modulo four is unique.

The norms and valuations giving independence are literal:
\[
 N(d)=27,\qquad
 N(\alpha)=1-\frac{8\cdot43^4}{27}\in1+8\Z_2,
 \qquad v_\theta(t_0)=1,\qquad
 v_\theta(\alpha-1)=3.
\tag{A.193.7}
\]
The second norm is a dyadic square and the first is not.
Moreover \(\alpha\) is not a square in \(L\).
Indeed every unit \(b\) has residue one.  Put
\(h=v_\theta(b-1)\geq1\).  If \(h<4\),
then \(v_\theta(b+1)=h\) and
\(v_\theta(b^2-1)=2h\in\{2,4,6\}\).
If \(h\geq4\), both factors have valuation at least four,
so \(v_\theta(b^2-1)\geq8\), with the zero case included.
None equals three.  Thus a product
\(d^a t_0^b\alpha^c\), with \(a,b,c\in\{0,1\}\),
can be a square only when \(b=0\) by valuation,
then \(a=0\) by its norm, and then \(c=0\) by this unit test.
The three classes in (A.193.6) are independent.

For completeness, the mod-two dimension is exactly three.
A sufficiently small formal subgroup of \(J(\Q_2)\)
is a rank-two \(\Z_2\)-module under the convergent logarithm.
It has finite index because \(J\) is proper.
For this subgroup multiplication by two has index four.
In the finite quotient the kernel and cokernel of multiplication
by two have equal orders; the multiplication diagram therefore
gives
\[
       |J(\Q_2)/2J(\Q_2)|=4\,|J(\Q_2)[2]|=8.
\tag{A.193.8}
\]
Here \(J(\Q_2)[2]=\{0,T\}\): the five finite branch points
have exactly the two Galois orbits \(\{0\}\) and the four roots
of the Eisenstein quartic.  The even-branch-divisor description
from A.186 gives precisely this one-dimensional invariant group.

The module \(\mathcal M\) is finitely generated over \(\Z_2\),
as follows from the same finite-index formal subgroup.
The three independent mod-two classes give, by the complete
Nakayama argument, a surjective continuous map
\[
 \Z_2\oplus\Z_2\oplus\Z/2\longrightarrow\mathcal M,
             \qquad (a,b,e)\longmapsto aD+bE_*+eT.
\tag{A.193.9}
\]
The map is injective as well.  Use the actual functional
\(\Lambda(L_1,L_2)=L_2-\rho L_1\) of A.189,
where \(\rho\equiv1\pmod4\).  It vanishes on \(D,T\),
while the positive-valuation calculation of A.192 with \(k=1\)
gives
\[
                   v_2\bigl(\Lambda(\log E_*)\bigr)=1.
\tag{A.193.10}
\]
The logarithm of \(D\) is nonzero.  Applying \(\Lambda\)
and then the logarithm to a kernel element in (A.193.9)
forces \(b=a=0\); the nonzero mod-two Kummer class of
\(T\) then forces \(e=0\).  This proves the actual direct
sum (A.193.2), including its topology, and also proves
\[
                  \Lambda(\log\mathcal M)=2\Z_2.
\tag{A.193.11}
\]
The logarithm on the completion is the continuous extension
of the local logarithm.  The kernel of
\(J(\Q_2)\to\mathcal M\) is its finite odd-order torsion
subgroup, which the logarithm kills, so this map is well defined.

\emph{Exact finite Selmer membership.}
The global completed group is the separate group
\(\Z_2D\oplus(\Z/2)T\), by the actual generator comparison
in A.186.  Since \(\Sha(J/\Q)[2^\infty]=0\), the finite
Kummer sequences identify its reduction modulo \(2^n\)
with the entire \(\operatorname{Sel}^{2^n}(J/\Q)\).
Writing the local completed class of \(j_C(R)\) uniquely as
\(aD+bE_*+eT\), its finite Kummer class is in the localized
global Selmer image exactly when \(b\in2^n\Z_2\).
By (A.193.10) this is equivalent to
\[
                 v_2\bigl(\Lambda(\log j_C(R))\bigr)
                                      \geq n+1.
\tag{A.193.12}
\]
For positive \(v_2(x)=2k-1\), A.192 gives the left-hand
valuation exactly \(k\).  For negative \(v_2(x)=-2k\)
it gives exactly \(k+1\).
For a unit \(x=3+h\), \(h\ne0\), set \(e=v_2(h)\geq3\).
Subtract the appropriate actual point \((3,\pm9)\), with
the same residue of \(y\) modulo four.  A.189's exact
integrals satisfy \(v_2(I_1)=e\) and
\(I_2/I_1\equiv3\pmod4\).  As \(\rho\equiv1\pmod4\),
their difference has valuation exactly one, and hence
\(v_2\Lambda=e+1\).  The four listed rational centres have
\(b=0\).  This proves every direction of (A.193.3).
Each fixed finite test still admits infinitely many ordinary
points close to \((3,9)\), since all \(x\in3+2^e\Z_2\)
with \(e\geq\max(3,n)\) have the required square root and
pass this test.  This assertion concerns the stated finite
Jacobian test, without asserting a locally curve-matching
global quotient section for each such point.

\emph{The actual finite arithmetic quotient.}
Let \(H\subset\Pi\) be the cyclic-eight cover subgroup from
A.192.  Compactification and the rational origin \(\infty\)
give the geometric map
\[
 H\longrightarrow\pi_1(C_{\overline\Q})^{\rm ab}
       \longrightarrow T_2J/2^nT_2J=J[2^n].
\tag{A.193.13}
\]
Let \(K_n\) be its kernel.  It is normal in the full arithmetic
tripod group: deck maps extend to the proper curve, permute its
boundary, and preserve the closed commutator and multiplication
subgroups.  Thus the actual quotient is
\(E_n=\pi_1(U)/K_n\), with geometric group \(Q_n=\Pi/K_n\).
It receives the full metabelian quotient because
\(\Pi''\subset\overline{[H,H]}\subset K_n\).

Here is an explicit integral presentation.  Take standard
geometric meridians \(x,y\), with characters one and three
modulo eight, and set \(v=x,u=yx^{-3}\).
The cyclic-cover Schreier basis consists of
\(u_j=v^juv^{-j}\), \(0\leq j<8\), and \(z=v^8\).
In its abelianization, the proper boundary relations are
\[
 z=0,\qquad u_j+u_{j+4}+z=0\ (0\leq j<4),
       \qquad\sum_{j=0}^7u_j+3z=0.
\tag{A.193.14}
\]
They come respectively from \(x^8\), the four lifts of
\((xy)^2\), and \(y^8\); the last relation is implied by
the first five.  Consequently the compactified integral
module is
\(\Z_2[v]/(v^4+1)=\Z_2[\zeta_8]\), generated by \(u_0\).
The boundary relation \(v^8=1\) gives a literal splitting
of the finite geometric extension.  The group in (A.193.4)
therefore has the explicit map
\[
                    x\longmapsto v,\qquad
                    y\longmapsto u_0v^3.
\tag{A.193.15}
\]
This also proves its order, without a polarization assumption.
Put \(\delta=\zeta_8-1\).  The polynomial
\((1+\delta)^4+1\) is Eisenstein at two, so its integral
ring is a discrete valuation ring with
\(v_\delta(2)=4\).  In the displayed semidirect product,
\(\Gamma_kQ_n=\delta^{k-1}(\Z_2[\zeta_8]/2^n)\)
for \(k\geq2\).  The first vanishing power is
\(\delta^{4n}\), whereas \(\delta^{4n-1}\ne0\).
This proves exact nilpotency class \(4n\).

A section in the stated combined class has a fixed sheet
of the actual eighth-root torsor.  Its image in the cover
subgroup of \(E_n\), followed by (A.193.13), is the proper
finite geometrically abelian extension split by \(\infty\).
The selected local point or cusp sheets therefore give exactly
the finite Kummer classes of their proper centres.  This is
the finite version of A.192's map, including the real component.
The resulting global class lies in
\(\operatorname{Sel}^{2^n}(J/\Q)\), proving the asserted
necessary finite test for quotient sections.

Finally \(11^5-54\cdot11=160457\equiv1\pmod8\), so the
two points with \(x=11\) are actual dyadic points of \(W\).
Their images are (A.193.5).  Every dyadic point of that
torsor fibre differs by \(\mu_8(\Q_2)=\{1,-1\}\), so the
inverse formulas in (A.192.12) preserve \(x=11\) and negate
\(y\).  Both have \(v_2(x-3)=3\).  They fail the test at
\(n=4\), giving the stated order-\(2^{19}\), class-sixteen
finite separator on every possible selected sheet.
\end{proof}
